\section{A budget turns subset search into DP}
\label{case:knapsack}
Synthetic task: each item has integer weight $w_i\ge0$ and value $v_i$.
Choose each item at most once, with total weight at most $C$, maximizing
value. There are $2^n$ subsets. When $nC$ fits the runtime and $C$ fits
memory, budget is a useful state dimension: all choices among the first $i$
items with the same remaining capacity have the same future options.

Define $D_i[c]$ as the best value using the first $i$ items with weight at
most $c$. Empty selection is allowed, so $D_0[c]=0$. Either skip item $i$,
or take it from the previous row:
\[
 D_{i+1}[c]=\max\bigl(D_i[c],D_i[c-w_i]+v_i\bigr)
 \quad(c\ge w_i).
\]
If it does not fit, only the skip term remains. This partitions every valid
subset according to whether it contains item $i$, proving the transition.
It costs $O(nC)$ rather than $O(2^n)$; huge binary-encoded capacities can
still make it impractical. If the sum of values is small instead, consider
minimum weight per value. A small number of items with huge weights may
instead suggest meet-in-the-middle enumeration.

Trace weights $[2,3,4]$, values $[4,5,7]$, capacity 5. The first row after
item 0 is $[0,0,4,4,4,4]$. After item 1 it is $[0,0,4,5,5,9]$; capacity 5
uses the previous value 4 at capacity 2 plus the new value 5. Item 2 improves
capacity 4 to 7 but does not beat 9 at capacity 5. The answer is 9, choosing
the first two items. The implementation book's
\verb@finals_core::knapsack01({2,3,4},{4,5,7},5)@ computes this value.

Only the previous row is needed for values. In a single array, scan
capacities downward so $D[c-w_i]$ still belongs to that previous row. For
one item $(w,v)=(2,3)$ and $C=4$, scanning upward writes $D[2]=3$, then
uses it again to write $D[4]=6$. That is two copies of one item and is wrong
for 0--1 selection. It is exactly the intended behavior for an unbounded
positive-weight item, where ascending order allows reuse.

If the total weight must equal $C$, initialize only $D[0]=0$ and all other
states as unreachable. Skip transitions out of unreachable states; an
at-most initialization would incorrectly treat unattainable weights as free.
To recover item indices, retain the rows and walk backward through skip/take
choices. The supplied one-row helper returns only the optimal value.

Practice: weights $[2,4]$, values $[3,8]$, exact weight 3. Answer: unreachable.
For weight at most 3 the answer is 3. This single change in the output
constraint changes the base cases, even though the recurrence looks similar.

\section{Turn a worked case into a new solution}
Use the examples as derivations, not text-matching rules. For a new task,
write down these differences before reusing code:
\begin{enumerate}
\item \textbf{Feasibility.} Are equal coordinates allowed? May vertices
repeat? Must every item be selected or assigned? Are numbers positive?
\item \textbf{Objective.} Existence, count, minimum cost, lexicographically
smallest answer, and reconstruction need different stored information.
\item \textbf{Scale.} Count the reduced states using the new bounds.
Remainder DP is useful only if length times modulus is affordable.
\item \textbf{Timing and order.} Are updates online or offline? Are equal
keys simultaneous? Can a hazard wait? Are dependencies acyclic?
\item \textbf{Output.} Does the helper return a count, parent indices, a
value, or an optional answer? Translate it into the statement's format.
\end{enumerate}

\subsection{How to move from the naive recurrence}
Start with a mathematically correct slow solution on small inputs. Then
look at the operation inside its innermost loop:
\begin{itemize}
\item A contiguous sum: prefix sums remove the inner loop.
\item Best earlier value with coordinate below $x$: process one coordinate
in order and query a prefix maximum in the other. Preserve strict ties.
\item Repeating the same suffix after different histories: memoize a state
that retains exactly the history affecting that suffix.
\item Rechecking every capacity: prove feasibility monotone, then search
its boundary. The feasibility test itself still needs to be efficient.
\item Trying all pairs of large objects: count equal features or
contributions. Do not assume grouping by one feature preserves every term.
\end{itemize}
Fundraising demonstrates the second change. Enigma demonstrates the third.
Displaying Decimals shows why deriving the value for one pair does not by
itself solve aggregation over all pairs.

\section{What to test before trusting the adaptation}
Choose tests that attack the reduction. A random test of the final numeric
answer does not explain why it is correct, but disagreement with a slow
oracle immediately disproves it.

For each example below, write an enumerator that directly follows the task's
rules, without reusing the optimized recurrence:
\begin{itemize}
\item Donation chains: enumerate subsets and check every selected pair's
compatibility. This catches illegal ties and duplicate handling.
\item Contiguous grouping: enumerate all cuts between adjacent jobs and
minimize the largest group sum. This catches an incorrect greedy predicate.
\item Toggle actions: BFS over board states with one click per edge. Compare
the shortest distance with action-subset enumeration.
\item Hazard escape: advance all possible monster positions and valid
player positions one time layer at a time. Compare existence and path length;
then replay the returned path for walls and unsafe simultaneous arrivals.
\item Exact sum: enumerate every pair of endpoints and add its contents.
This catches prefix initialization and query-before-insert mistakes.
\end{itemize}
These are the independent oracles used for this edition's five new modules.
They exercise small cases; numeric bounds and maximum-size complexity still
need separate checks. Named contest snippets have local test evidence, not
new online-judge verdicts.

\subsection{Counterexamples worth writing by hand}
Equal keys; all values zero; one item; no solution; the start already at the
goal; disconnected components; two different histories reaching the same
vertex; an exact capacity boundary; a prefix sum repeated several times;
a valid answer equal to the chosen sentinel. For each, predict the result
before running the code. If the result is ambiguous, the model or output
contract is not yet precise enough.

\section{Practice without choosing another topic}
Take one case at a time. Read its recognition and reduced task, cover the
implementation, and write the state or invariant. Trace the given example.
Implement it, compare small cases with the oracle, then answer the variation
that changes one assumption. Finish by reopening a past statement without
its algorithm name and explaining why the same method does or does not fit.

A concrete order for the team's reported gaps:
\begin{enumerate}
\item Signed prefix counts and sorted interval union. Make endpoints,
strictness, and multiplicity reliable.
\item Parity-state BFS, then the hazard deadline model. Draw states and
transitions before writing queue code.
\item Enigma. Derive the state, backward feasibility, and forward smallest
reconstruction; do not merely copy the three loops.
\item Hidden Number and capacity search. Give the greedy proof in words,
then deliberately change positivity or ordering and see what fails.
\item Fundraising. Implement the quadratic recurrence first, then replace
its predecessor scan with a range maximum and test tied coordinates.
\item Fuzzy Factorization and single-fraction decimals. Separate exact
representation from approximation and prove the chosen error convention.
\item One assignment/flow example. Name vertices, edges, capacities, and
what one unit of flow represents before choosing the routine.
\end{enumerate}

\section{When the current pattern is insufficient}
The original KACTL has additional specialists. Choose them only after the
model demands their extra capability. The following are navigation notes,
not complete implementations in this workbook:
\begin{itemize}
\item Dynamic tree path aggregates: heavy-light decomposition plus a range
structure. Edge versus vertex values and update composition matter.
\item Offline queries with reversible local changes: Mo ordering. You must
supply correct add/remove operations; it is not a ready-made answer.
\item A recurrence that sums products across all offsets: convolution via
FFT/NTT. Check transform length, modulus, and coefficient bounds first.
\item Many substring-order queries: suffix array/LCP. KMP and Z solve
prefix/matching tasks, not arbitrary suffix ordering.
\item Weighted assignment: min-cost flow or an assignment algorithm.
Cardinality matching ignores weights; ordinary max flow ignores costs.
\item Large integer factorization: algorithms beyond trial division.
A factorization routine still does not replace a reduction that avoids
factoring the original giant number.
\end{itemize}
A novel problem may combine several models or require a new proof. When a
case's assumptions do not fit, return to the naive formulation and find the
missing state or operation. Adding more templates without identifying that
mismatch will not repair the solution.
